\medskip








\begin{enumerate}
\item 
\[f(x) = x \e^{x^2 - 3}.\]








Avec $F(x) = \dfrac12 \e^{x^2 - 3}$, on a $F'(x) = 2x \times \dfrac12\e^{x^2 - 3} = x\e^{x^2 - 3} = f(x)$ : réponse \textbf{d.~}






\item On considère la suite $\left(u_n\right)$ définie pour tout entier naturel $n$ par: \[u_n = \e^{2n+1}.\]









Si pour tout entier naturel $n$, \: $u_n = \e^{2n+1}$, alors $u_{n+1} = \e^{2(n+1)+ 1} = \e^{2n + 2 + 1} = \e^{2} \times \e^{2n + 1} = \e^{2} \times u_n$ : cette égalité montre que la suite $\left(u_n\right)$ est géométrique de raison $\e^{2}$ : réponse \textbf{c.~~}

Pour les questions \textbf{3.} et \textbf{4.}, on considère la suite $\left(u_n\right)$ définie sur $\N$ par:
\[u_0 = 15 \quad \text{et pour tout entier naturel }\: n \quad :\: u_{n+1} = 1,2u_n + 12.\]

\item ~

\begin{center}
\begin{tabular}{|>{\ttfamily}l|}\hline
def seuil() :\\
\qquad n=0\\
\qquad u=15\\
\qquad while u$\leqslant \np{10000}$:\\
\qquad \quad n=n+1\\
\qquad \quad u=1,2$*$u+12\\
\qquad return(n)\\ \hline
\end{tabular}
\end{center}

Réponse \textbf{a.}








\item On considère la suite $\left(v_n\right)$ définie sur $\N$ par : $v_n = u_n + 60$. 









On a quel que soit $n \in \N, \: v_{n+1} = u_{n+1} + 60 =  1,2u_n + 12 + 60 = 1,2u_n + 72 = 1,2\left(u_n + 60\right) = 1,2v_n$ : cette égalité montre que la suite $\left(v_n\right)$ est géométrique de raison 1,2.
\end{enumerate}

\bigskip
